% =============================================================================
% References
% =============================================================================
\section{References}\label{app:refs}

The first edition of this book cited its sources author-and-year in prose and
supplied no list, which left a reader no way to check them. That is a poor
arrangement in any textbook and an indefensible one here, since the book's
argument is that a claim should be traceable to the measurement behind it. The
list below closes the gap.

Each entry carries a DOI where we could verify one against a primary record, and
says so where we could not --- three entries are marked \textbf{[no verified
DOI]} rather than given a plausible-looking string, which for a pre-1970 paper
is the honest outcome and for this book the only tolerable one. The annotation
says what the source is used \emph{for}, because a reference list that does not
is only a bibliography of things the authors have heard of. Where this book
relies on a specific table or figure, the location is named; the machine-readable
version, with verbatim quotations of the sentences the numbers come from, is
\file{data/pole_proofreading_literature.json}.

\subsection*{The organising case}

\begin{refs}
\item \textbf{Kane DP, Shcherbakova PV} (2014). A common cancer-associated DNA
polymerase $\varepsilon$ mutation causes an exceptionally strong mutator
phenotype, indicating fidelity defects distinct from loss of proofreading.
\emph{Cancer Research} \textbf{74}(7):1895--1901.
\doi{10.1158/0008-5472.CAN-13-2892}
\\[0.15em]
\emph{Used in:} §\ref{sec:refutation}, Table~\ref{tab:kane}. The source of the
three \texttt{CAN1} rates --- wild type $15$, \texttt{pol2-4} $44$,
\texttt{pol2-P301R} $2300$, all $\times10^{-8}$ with $95\%$ CIs --- whose ratio
$\rho = 52.3$ refutes the stage-removal class. The title states the paper's own
conclusion, reached independently of any theorem in this book.

\item \textbf{Xing X, Kane DP, Bulock CR, Moore EA, Sharma S, Chabes A,
Shcherbakova PV} (2019). A recurrent cancer-associated substitution in DNA
polymerase $\varepsilon$ produces a hyperactive enzyme. \emph{Nature
Communications} \textbf{10}:374. \doi{10.1038/s41467-018-08145-2}
\\[0.15em]
\emph{Used in:} §\ref{sec:refutation}, §\ref{sec:mmr}. The purified-enzyme
comparison ($59$ against $155 \times 10^{-6}$ base substitutions) that creates
the reversal Section~\ref{sec:mmr} is built on, and the tetrad analysis
establishing that \texttt{pol2-P301R msh6$\Delta$} is synthetic lethal. Its
central result --- that the variant is a hyperactive polymerase --- is also the
reason Remark~\ref{rem:transport} declines to call this book's capacity
conclusion unconditionally forced.

\item \textbf{Barbari SR, Kane DP, Moore EA, Shcherbakova PV} (2018). Functional
analysis of cancer-associated DNA polymerase $\varepsilon$ variants in
\emph{Saccharomyces cerevisiae}. \emph{G3: Genes, Genomes, Genetics}
\textbf{8}(3):1019--1029. \doi{10.1534/g3.118.200042}
\\[0.15em]
\emph{Used in:} §\ref{sec:refutation}, control~1. Thirteen further
cancer-associated alleles, of which three exceed the PSR$_2$ ceiling. This paper
has no \texttt{pol2-4} arm, which is why those three ratios require a cross-paper
normalisation and are flagged as weaker than the within-table comparison.

\item \textbf{Barbari SR, Shcherbakova PV}, \emph{et al.}\ (2022). Enhanced
polymerase activity permits efficient synthesis by cancer-associated DNA
polymerase $\varepsilon$ variants at low dNTP levels. \emph{Nucleic Acids
Research} \textbf{50}(14):8023--8040. \doi{10.1093/nar/gkac602}
\\[0.15em]
\emph{Used in:} §\ref{sec:refutation} control~2, Worked Example~7.1. The source
that overturned this book's description of \texttt{V411L} as functionally
neutral: the analogous yeast allele is a robust mutator in an MMR-deficient
background and the purified enzyme is among the most active tested.

\item \textbf{Bulock CR, Xing X, Shcherbakova PV} (2020). Mismatch repair and DNA
polymerase $\delta$ proofreading prevent catastrophic accumulation of leading
strand errors in cells expressing a cancer-associated DNA polymerase
$\varepsilon$ variant. \emph{Nucleic Acids Research} \textbf{48}(16):9124--9134.
\doi{10.1093/nar/gkaa633}
\\[0.15em]
\emph{Used in:} §\ref{sec:mmr}. Establishes that two distinct repair routes are
required for viability, which is why the capacity $c$ of
Assumption~\ref{model:channel} must be read as lumping them together.

\item \textbf{Parkash V}, \emph{et al.}\ (2019). Structural consequence of the most
frequently recurring cancer-associated substitution in DNA polymerase
$\varepsilon$. \emph{Nature Communications} \textbf{10}:373.
\doi{10.1038/s41467-018-08114-9}
\\[0.15em]
\emph{Used in:} not relied upon for any number in this book. Listed because it is
the companion crystal structure to Xing \emph{et al.}\ (2019) and is the natural
next reading for a student who wants the structural account of why the arginine
substitution behaves as it does.
\end{refs}

\subsection*{Kinetic proofreading}

\begin{refs}
\item \textbf{Hopfield JJ} (1974). Kinetic proofreading: a new mechanism for
reducing errors in biosynthetic processes requiring high specificity.
\emph{Proceedings of the National Academy of Sciences USA}
\textbf{71}(10):4135--4139. \doi{10.1073/pnas.71.10.4135}
\\[0.15em]
\emph{Used in:} §\ref{sec:proofreading}. The original argument that irreversible
delay stages purchase discrimination beyond the equilibrium limit.

\item \textbf{Ninio J} (1975). Kinetic amplification of enzyme discrimination.
\emph{Biochimie} \textbf{57}(5):587--595.
\doi{10.1016/S0300-9084(75)80139-8}
\\[0.15em]
\emph{Used in:} §\ref{sec:proofreading}. Independent and near-simultaneous; the
mechanism is properly credited to both.
\end{refs}

\subsection*{Mutation rate estimation}

\begin{refs}
\item \textbf{Luria SE, Delbrück M} (1943). Mutations of bacteria from virus
sensitivity to virus resistance. \emph{Genetics} \textbf{28}(6):491--511.
\doi{10.1093/genetics/28.6.491}
\\[0.15em]
\emph{Used in:} §\ref{sec:accumulation}. The fluctuation test, and the argument
that the variance rather than the mean distinguishes pre-existing mutation from
induced adaptation.
\end{refs}

\subsection*{The structure of DNA}

\begin{refs}
\item \textbf{Cochran W, Crick FHC, Vand V} (1952). The structure of synthetic
polypeptides. I. The transform of atoms on a helix. \emph{Acta
Crystallographica} \textbf{5}:581--586. \nodoi
\\[0.15em]
\emph{Used in:} §\ref{sub:photo51}, equations \eqref{eq:ccv} and
\eqref{eq:interf}. The helical diffraction theory from which
Figure~\ref{fig:photo51} is recomputed, including the interference term that
makes the fourth layer line vanish at $\delta = 3/8$.

\item \textbf{Franklin R, Gosling RG} (1953). Molecular configuration in sodium
thymonucleate. \emph{Nature} \textbf{171}(4356):740--741.
\doi{10.1038/171740a0}
\\[0.15em]
\emph{Used in:} the dedication and §\ref{sub:photo51}. The published account of
the B-form fibre diffraction data, including the helical parameters. Photograph
51 itself was exposed in May 1952.

\item \textbf{Watson JD, Crick FHC} (1953). Molecular structure of nucleic acids:
a structure for deoxyribose nucleic acid. \emph{Nature}
\textbf{171}(4356):737--738. \doi{10.1038/171737a0}
\\[0.15em]
\emph{Used in:} the dedication. The model that the data above constrained, and
whose acknowledgements did not say so.
\end{refs}

\subsection*{DNA topology}

\begin{refs}
\item \textbf{Călugăreanu G} (1961). Sur les classes d'isotopie des n\oe{}uds
tridimensionnels et leurs invariants. \emph{Czechoslovak Mathematical Journal}
\textbf{11}(4):588--625. \doi{10.21136/CMJ.1961.100486}
\\[0.15em]
\emph{Used in:} §\ref{sec:dnatop}. The first form of the $\Lk = \Tw + \Wr$
relation.

\item \textbf{White JH} (1969). Self-linking and the Gauss integral in higher
dimensions. \emph{American Journal of Mathematics} \textbf{91}(3):693--728.
\nodoi{} JSTOR \texttt{2373348}.
\\[0.15em]
\emph{Used in:} §\ref{sec:dnatop}, Theorem~\ref{thm:cwf}. The general proof,
without the curvature restriction. This book cites the theorem and does not
prove or formalise it; only its consequence, Theorem~\ref{thm:trade}, is
machine-checked.

\item \textbf{Fuller FB} (1971). The writhing number of a space curve.
\emph{Proceedings of the National Academy of Sciences USA}
\textbf{68}(4):815--819. \doi{10.1073/pnas.68.4.815}
\\[0.15em]
\emph{Used in:} §\ref{sec:dnatop}, Definition~\ref{def:lktwwr}. The writhe, and
the name.
\end{refs}

\subsection*{Methods and software}

\begin{refs}
\item \textbf{The GUDHI Project} (2015--). \emph{GUDHI User and Reference Manual}.
INRIA. \url{https://gudhi.inria.fr/}
\\[0.15em]
\emph{Used in:} §\ref{sec:tda}. Version $3.13.0$; Vietoris--Rips persistence on
the chaos-game point cloud.

\item \textbf{Chan JM, Carlsson G, Rabadan R} (2013). Topology of viral
evolution. \emph{Proceedings of the National Academy of Sciences USA}
\textbf{110}(46):18566--18571. \doi{10.1073/pnas.1313480110}
\\[0.15em]
\emph{Used in:} Appendix~\ref{app:directions}, direction D3. The construction
that makes persistent homology informative on sequence data, and the one this
book should have used in Section~\ref{sec:tda}: compute persistence on a
\emph{population} under a genetic-distance metric, not on an embedding of a
single sequence. Clonal evolution is summarised by a tree; reticulate evolution
exhibits nontrivial topology above dimension zero. The contrast with
\S\ref{sub:tdanull} is the useful part --- there a loop had no meaning, so the
null had to be an ensemble of shuffles; here a loop means non-tree-like
evolution, and the null can be a theorem.

\item \textbf{Bauer U, Roll F} (2021). Gromov hyperbolicity, geodesic defect,
and apparent pairs in Vietoris--Rips filtrations. \texttt{arXiv:2112.06781}.
\nodoi
\\[0.15em]
\emph{Used in:} Appendix~\ref{app:directions}, direction D3. The source for the
claim that finite tree metrics have Vietoris--Rips complexes collapsing onto
their subforests, hence acyclic. We cite this for the tree-metric case and
prove only the sharper ultrametric case ourselves, in
\file{formal/ANSE/ReticulationCertificate.lean} --- having searched for a
citable statement of it and not found one. The asymmetry is recorded in the code,
which labels four-point defects \emph{suggestive} and ultrametric certificates
\emph{proved}.

\item \textbf{Morozov D, Nigmetov A} (2020). Brief announcement: towards
lockfree persistent homology. In \emph{Proceedings of the 32nd ACM Symposium on
Parallelism in Algorithms and Architectures}, 555--557.
\doi{10.1145/3350755.3400244}
\\[0.15em]
\emph{Used in:} the Rust implementation of direction D3, via
\href{https://github.com/tomchaplin/lophat}{LoPHAT}, which implements this
algorithm. Our crate supplies the Vietoris--Rips front-end that LoPHAT
deliberately omits.

\item \textbf{Altschul SF, Erickson BW} (1985). Significance of nucleotide
sequence alignments: a method for random sequence permutation that preserves
dinucleotide and codon usage. \emph{Molecular Biology and Evolution}
\textbf{2}(6):526--538. \doi{10.1093/oxfordjournals.molbev.a040370}
\\[0.15em]
\emph{Used in:} \S\ref{sub:tdanull}. The Eulerian-path shuffle that preserves
every dinucleotide count exactly, which is the \textbf{dinuc} null of
Assumption~\ref{model:tdanull}. The paper also supplies the codon-preserving
permutation used as the third null, and makes the same methodological point this
book makes in Section~\ref{sec:accumulation}: that a shuffle is only a null if
you can say exactly what it preserves.

\item \textbf{Gotoh O} (1982). An improved algorithm for matching biological
sequences. \emph{Journal of Molecular Biology} \textbf{162}(3):705--708.
\doi{10.1016/0022-2836(82)90398-9}
\\[0.15em]
\emph{Used in:} \S\ref{sub:ortho}. The affine-gap alignment recursion
implemented in
\file{scripts/textbook/verify_yeast_human_correspondence.py}. Worth knowing: the
1982 paper is reported to contain indexing and matrix-initialisation errors that
can yield sub-optimal alignments, so an implementation should not be trusted
because it follows the paper. Ours is checked against three cases whose answers
are known independently --- a self-alignment must be the identity map, a
single-residue deletion must place exactly one gap at that residue, and affine
costs must prefer one long gap to several short ones --- in
\file{tests/test_tda_null_and_orthologue.py}.

\item \textbf{Jeffrey HJ} (1990). Chaos game representation of gene structure.
\emph{Nucleic Acids Research} \textbf{18}(8):2163--2170.
\doi{10.1093/nar/18.8.2163}
\\[0.15em]
\emph{Used in:} §\ref{sec:tda}, Definition~\ref{def:cgr}. The embedding, in its
original two-dimensional form; this book uses the tetrahedral extension to
$\mathbb{R}^3$.

\item \textbf{de Moura L, Kong S, Avigad J, van Doorn F, von Raumer J} (2015).
The Lean theorem prover (system description). In \emph{Automated Deduction ---
CADE-25}, Lecture Notes in Computer Science 9195, 378--388. Springer. \nodoi
\\[0.15em]
\emph{Used in:} throughout. The proofs in this book are checked with Lean~4 and
\emph{mathlib}; the gate described in §\ref{sec:proof-and-cell} adds the
\lean{\#print axioms} audit on top of compilation.
\end{refs}

\subsection*{A note on what is not cited here}

Three absences are deliberate.

The pembrolizumab TMB threshold in §\ref{sec:clinic} is stated as a regulatory
criterion rather than attributed to a trial, because the book makes no claim
about the evidence base behind it beyond the fact that the decision rule takes
measured burden as an input.

The $\Delta\Delta G \approx 4.6\,k_BT$ of §\ref{sec:proofreading} is a standard
thermodynamic estimate rather than a specific measurement, and is used only to
establish an order of magnitude.

Earlier editions of this material carried a reference, \texttt{anthropic2026art},
that did not correspond to a real publication. It was removed when discovered.
The incident is the reason every entry above was re-verified against a primary
record, with DOIs, rather than carried forward from a previous draft.
